\begin{pblproject}[Checkpoint project: path length, and composition respects it]
\label{sec:11-path-algebras:07-checkpoint-project:1}


Part II closes here, having built groups, rings, modules, and now quivers and paths from scratch. This project extends the \texttt{Path}/\texttt{Path.append} construction from §4--§5 with one more piece of structure --- a path's \textbf{length} --- and proves it behaves the way it obviously should under composition, tying together this chapter's inductive-type work with the ``prove it once, generically'' habit Chapters 7 and 9 established.

\textbf{Learning objectives.} Define a recursive function over the same indexed inductive type \texttt{Path} (§4), then prove a genuine theorem about it by induction, mirroring \texttt{Path.append}'s own recursion case for case (§5). Along the way, this project surfaces a real, verifiable fact about Lean: because \texttt{Path} is \emph{indexed} by both its endpoints, functions matching on it (both \texttt{Path.append} and the \texttt{Path.length} built here) do not reduce by bare \texttt{rfl} once an abstract path variable is involved --- only their auto-generated equation lemmas do. This is worth discovering directly rather than being told, the same way §5's own worked proofs are best understood by predicting each \texttt{rw}'s effect before running it.

\textbf{Prerequisites.} Chapters 1--11, specifically §4 (\texttt{Path}) and §5 (\texttt{Path.append}).

\textbf{Milestones.}

\begin{enumerate}
\def\labelenumi{\arabic{enumi}.}
\tightlist
\item
  Define \texttt{Path.length : \textbackslash{}\{u v : V\textbackslash{}\} → Path Q u v → Nat} by recursion: the trivial path \texttt{nil} has length \texttt{0}; \texttt{cons a h h' p} has length \texttt{p.length + 1} (one more than the path it extends).
\item
  Compute a couple of lengths by hand and check them with \texttt{\textbackslash{}\#eval} against §4's example paths (\texttt{pathAlpha}, \texttt{pathBetaAlpha}).
\item
  Prove \texttt{Path.append\textbackslash{}\_length : (Path.append p q).length = p.length + q.length} by induction on \texttt{q}, mirroring \texttt{Path.append}'s own \texttt{nil}/\texttt{cons} cases. Predict, before running it, whether each case closes by \texttt{rfl} --- then discover (as the note above flags) that it does not, and that \texttt{simp only [Path.append, Path.length]} is needed to unfold one step, exactly as \texttt{rfl} needs the \emph{equation lemmas} Lean generates for an indexed match, not raw iota-reduction, once an abstract path is involved.
\end{enumerate}

\textbf{Deliverable.} \texttt{Path.length} and the proved theorem \texttt{Path.append\textbackslash{}\_length}, checked against at least one concrete instance.

\textbf{Self-verification.}

\begin{lstlisting}
def Path.length {V A : Type} {Q : Quiver V A} : {u v : V} → Path Q u v → Nat
  | _, _, Path.nil _ => 0
  | _, _, Path.cons _ _ _ p => p.length + 1

#eval pathAlpha.length        -- 1
#eval pathBetaAlpha.length    -- 2

theorem Path.append_length {V A : Type} {Q : Quiver V A} {u v w : V}
    (p : Path Q u v) (q : Path Q v w) :
    (Path.append p q).length = p.length + q.length := by
  induction q with
  | nil =>
    simp only [Path.append, Path.length]
    rw [Nat.add_zero]
  | cons a h h' q' ih =>
    simp only [Path.append, Path.length]
    rw [ih, Nat.add_assoc]

-- The concrete check: pathBetaAlpha was §5's own worked example of
-- Path.append; its length should be pathAlpha.length + pathBetaOnly.length.
example : (Path.append pathAlpha pathBetaOnly).length =
    pathAlpha.length + pathBetaOnly.length :=
  Path.append_length pathAlpha pathBetaOnly

#eval (Path.append pathAlpha pathBetaOnly).length   -- 2
\end{lstlisting}

If this compiles and the \texttt{\textbackslash{}\#eval}s match, the project is done. A full worked solution is in

\end{pblproject}
